% ============================================================
%  厦门大学本科毕业论文（设计）LaTeX 模板
%  编译方式：XeLaTeX
% ============================================================
\documentclass[a4paper,12pt,openright]{book}

% ==================== 基础宏包 ====================
\usepackage{xeCJK}
\usepackage{ctex}
\usepackage{fontspec}
\usepackage{amsmath,amsthm,amssymb,mathrsfs,bm}
\usepackage{latexsym}
\usepackage{float}
\usepackage{graphicx}
\usepackage{subfigure}
\usepackage{array}
\usepackage{tabularx}
\usepackage{multirow}
\usepackage{multicol}
\usepackage{booktabs}
\usepackage{caption}
\usepackage{fancyhdr}
\usepackage{titlesec}
\usepackage{titletoc}
\usepackage{hyperref}
\usepackage{url}
\usepackage{setspace}
\usepackage{appendix}
\usepackage{tikz}
\usepackage{listings}
\usepackage{xcolor}
\usepackage{CJKnumb}
\usepackage[final]{pdfpages}
\usepackage{verbatim}
\usepackage{algorithm}
\usepackage{algpseudocode}
\usepackage{algorithmicx}
\usepackage{bbm}

% ==================== 页面设置 ====================
% 页边距：上、左 ≥ 25mm，下、右 ≥ 20mm
\usepackage[top=25mm,bottom=20mm,left=25mm,right=20mm]{geometry}

% ==================== 字体设置 ====================
% 使用方式：
%   方案 A（Overleaf / 无系统字体）：上传字体到 fonts/ 文件夹，使用 Path 加载
%   方案 B（本地编译 + 已安装字体）：注释方案 A，取消方案 B 注释
%
% -------- 方案 A：Overleaf（通过文件路径加载上传的字体）--------
% 英文字体（Overleaf 自带 TeX Gyre 系列，与 Times New Roman / Arial 等价）
\setmainfont{texgyretermes}[
    Extension      = .otf,
    UprightFont    = *-regular,
    BoldFont       = *-bold,
    ItalicFont     = *-italic,
    BoldItalicFont = *-bolditalic,
]
\setsansfont{texgyreheros}[
    Extension      = .otf,
    UprightFont    = *-regular,
    BoldFont       = *-bold,
    ItalicFont     = *-italic,
    BoldItalicFont = *-bolditalic,
]
% 中文字体（Fandol 系列，Overleaf/TeX Live 自带，与宋黑楷仿等价）
\setCJKmainfont{FandolSong}[
    Extension   = .otf,
    UprightFont = *-Regular,
    BoldFont    = *-Bold,
]
\setCJKsansfont{FandolHei}[
    Extension   = .otf,
    UprightFont = *-Regular,
    BoldFont    = *-Bold,
]
\setCJKmonofont{FandolFang}[
    Extension   = .otf,
    UprightFont = *-Regular,
]
\setCJKfamilyfont{zhsong}{FandolSong}[Extension=.otf,UprightFont=*-Regular,BoldFont=*-Bold]
\setCJKfamilyfont{zhhei}{FandolHei}[Extension=.otf,UprightFont=*-Regular,BoldFont=*-Bold]
\setCJKfamilyfont{zhkai}{FandolKai}[Extension=.otf,UprightFont=*-Regular]
\setCJKfamilyfont{zhfs}{FandolFang}[Extension=.otf,UprightFont=*-Regular]

% -------- 方案 B：本地编译（macOS 已安装字体）--------
% 取消以下注释，并注释掉上方方案 A 的全部字体设置
% \setmainfont{Times New Roman}
% \setsansfont{Arial}
% \setCJKmainfont{Songti SC}[BoldFont=Songti SC Bold]
% \setCJKsansfont{PingFang SC}
% \setCJKmonofont{STFangsong}
% \setCJKfamilyfont{zhsong}{Songti SC}[BoldFont=Songti SC Bold]
% \setCJKfamilyfont{zhhei}{PingFang SC}
% \setCJKfamilyfont{zhkai}{STKaiti}
% \setCJKfamilyfont{zhfs}{STFangsong}

\newcommand*{\song}{\CJKfamily{zhsong}}  % 宋体
\newcommand*{\hei}{\CJKfamily{zhhei}}    % 黑体
\newcommand*{\kai}{\CJKfamily{zhkai}}    % 楷体
\newcommand*{\fs}{\CJKfamily{zhfs}}      % 仿宋

% ==================== 字号定义 ====================
\newcommand{\yihao}{\fontsize{26pt}{36pt}\selectfont}
\newcommand{\erhao}{\fontsize{22pt}{28pt}\selectfont}
\newcommand{\xiaoer}{\fontsize{18pt}{25pt}\selectfont}
\newcommand{\sanhao}{\fontsize{16pt}{24pt}\selectfont}
\newcommand{\xiaosan}{\fontsize{15pt}{22pt}\selectfont}
\newcommand{\sihao}{\fontsize{14pt}{21pt}\selectfont}
\newcommand{\banxiaosi}{\fontsize{13pt}{19.5pt}\selectfont}
\newcommand{\xiaosi}{\fontsize{12pt}{18pt}\selectfont}
\newcommand{\dawuhao}{\fontsize{11pt}{16.5pt}\selectfont}
\newcommand{\wuhao}{\fontsize{10.5pt}{15.75pt}\selectfont}
\newcommand{\xiaowu}{\fontsize{9pt}{13.5pt}\selectfont}

% ==================== 行距设置 ====================
% 正文 1.5 倍行距
\linespread{1.5}

% ==================== 图表标题设置 ====================
\renewcommand{\figurename}{图}
\renewcommand{\tablename}{表}
\captionsetup[table]{font={bf, small}, labelsep=quad}
\captionsetup[figure]{font={bf, small}, labelsep=quad}

% ==================== 代码环境设置 ====================
\lstset{
    numbers=left,
    numberstyle=\small\ttfamily,
    basicstyle=\small\ttfamily,
    keywordstyle=\small\ttfamily\color{blue},
    identifierstyle=\small\ttfamily,
    stringstyle=\small\ttfamily,
    commentstyle=\color[cmyk]{1,0,1,0},
    breaklines=true,
    extendedchars=false,
    showstringspaces=false,
    tabsize=4,
}

% ==================== 算法环境 ====================
\renewcommand{\algorithmicrequire}{\textbf{Input:}}
\renewcommand{\algorithmicensure}{\textbf{Output:}}

% ==================== 章节标题格式 ====================
% 章标题：小三号黑体，居中
\titleformat{\chapter}{\centering\xiaosan\hei\bfseries}{第\CJKnumber{\thechapter}章}{1em}{}
% 节标题：四号黑体
\titleformat*{\section}{\sihao\hei\bfseries}
% 小节标题：小四号黑体
\titleformat*{\subsection}{\xiaosi\hei\bfseries}

% 章标题占 2 行
\titlespacing{\chapter}{0pt}{0pt}{20pt}

% ==================== 目录格式 ====================
% 中文目录：章 - 四号黑体，节 - 小四号黑体，小节 - 小四号宋体
\titlecontents{chapter}[0pt]
{\addvspace{2pt}\filright\bfseries\sihao\hei}
{\contentspush{\thecontentslabel\hspace{1em}}}{}{~\titlerule*[6pt]{.}\contentspage}
\titlecontents{section}[20pt]
{\bfseries\xiaosi\hei}
{\contentspush{\thecontentslabel\hspace{1em}}}{}{~\titlerule*[6pt]{.}\contentspage}
\titlecontents{subsection}[40pt]
{\xiaosi\song}
{\contentspush{\thecontentslabel\hspace{1em}}}{}{~\titlerule*[6pt]{.}\contentspage}

% ==================== 定理环境 ====================
\newtheorem{thm}{定理}[section]
\newtheorem*{thm*}{定理}
\newtheorem{lmm}[thm]{引理}
\newtheorem{coro}[thm]{推论}
\newtheorem{prps}[thm]{命题}
\theoremstyle{definition}
\newtheorem{exmp}[thm]{例}
\newtheorem{dfn}[thm]{定义}
\theoremstyle{remark}
\newtheorem*{rmk}{注记}

% ==================== 页眉页脚设置 ====================
% 奇数页页眉：当前章名；偶数页页眉：论文题目
\pagestyle{fancy}
\renewcommand{\chaptermark}[1]{\markboth{第\CJKnumber{\thechapter}章~~~~#1}{}}
\fancyhead{}
\fancyhead[CO]{\xiaowu\song \leftmark}
%%%%%%%%%%%%%%%%%%%%%%%%%%%%%%%%%%%%%%%%%
%%%%% 请在下方修改论文标题 %%%%%
\fancyhead[CE]{\xiaowu\song 论文题目}
%%%%% 请在上方修改论文标题 %%%%%
%%%%%%%%%%%%%%%%%%%%%%%%%%%%%%%%%%%%%%%%%
\fancyfoot[CO,CE]{\xiaowu \thepage}

\fancypagestyle{plain}{
    \fancyhead{}
    \fancyfoot[CO,CE]{\xiaowu \thepage}
    \renewcommand{\headrulewidth}{0pt}
}

% ==================== 超链接设置 ====================
\hypersetup{
    colorlinks=true,
    linkcolor=black,
    citecolor=black,
    urlcolor=blue,
}

% ==================== 英文目录定义 ====================
\makeatletter
\newcommand\engcontentsname{Contents}
\newcommand\tableofengcontents{%
    \if@twocolumn
    \@restonecoltrue\onecolumn
    \else
    \@restonecolfalse
    \fi
    \chapter*{\engcontentsname
        \@mkboth{%
            \MakeUppercase\engcontentsname}{\MakeUppercase\engcontentsname}}%
    \@starttoc{toe}%
    \if@restonecol\twocolumn\fi
}
\newcommand\addengcontents[2]{%
    \addcontentsline{toe}{#1}{\protect\numberline{\csname the#1\endcsname}#2}}
\makeatother

\newcommand\echapter[1]{\addengcontents{chapter}{#1}}
\newcommand\esection[1]{\addengcontents{section}{#1}}
\newcommand\esubsection[1]{\addengcontents{subsection}{#1}}

% ==================== 数学快捷命令 ====================
\newcommand{\e}{\mathrm{e}}
\renewcommand{\i}{\mathrm{i}}
\renewcommand{\d}{\mathrm{d}}


% ============================================================
%                       正文开始
% ============================================================
\begin{document}

\frontmatter
\pagenumbering{Roman}

% ==================== 封面 ====================
% 方式一：引入外部封面 PDF
% \includepdf{titlepage.pdf}
% \newpage
% 方式二：使用 LaTeX 生成封面（取消下方注释使用）
\begin{titlepage}
    \begin{center}
        \vspace*{20mm}

        % 校名图片（需自行准备 xmu-logo.png）
        % \includegraphics[width=0.6\textwidth]{xmu-logo.png}\\[10mm]
        {\hei\yihao 厦\quad 门\quad 大\quad 学}\\[15mm]

        {\hei\sanhao 本\quad 科\quad 毕\quad 业\quad 论\quad 文\quad （设\quad 计）}\\[5mm]
        {\xiaosi （主修 / 辅修专业）}\\[15mm]

        % 中文标题：二号黑体
        {\hei\erhao 论文中文标题}\\[10mm]
        % 英文标题：三号 Times New Roman 加粗
        {\bfseries\sanhao English Title of the Thesis}\\[20mm]

        {\xiaosi
        \renewcommand\arraystretch{1.8}
        \begin{tabular}{r@{：}l}
            姓\hspace{2em}名 & 你的姓名 \\
            学\hspace{2em}号 & 你的学号 \\
            学\hspace{2em}院 & 你的学院 \\
            专\hspace{2em}业 & 你的专业 \\
            年\hspace{2em}级 & 你的年级 \\
            校内指导教师 & 教师姓名\hspace{1em}职称 \\
            校外指导教师 & 教师姓名\hspace{1em}职务 \\
        \end{tabular}
        }\\[20mm]

        {\xiaosi 二〇二六\quad 年\quad 月\quad 日}
    \end{center}
\end{titlepage}

\clearpage{\pagestyle{empty}\cleardoublepage}

% ==================== 诚信承诺书 ====================
% 方式一：引入外部 PDF
% \includepdf{Promise.pdf}
% \newpage
% 方式二：使用 LaTeX 生成
\chapter*{}
\thispagestyle{empty}
\begin{center}
    {\hei\xiaosan 厦门大学本科学位论文诚信承诺书}
\end{center}
\vspace{10mm}

{\xiaosi\song
本人呈交的学位论文是在导师指导下独立完成的研究成果。本人在论文写作中参考其他个人或集体已经发表的研究成果，均在文中以适当方式明确标明，并符合相关法律规范及《厦门大学本科毕业论文（设计）规范》。

该学位论文为（\hspace{6em}）课题（组）的研究成果，获得（\hspace{6em}）课题（组）经费或实验室的资助，在（\hspace{6em}）实验室完成（请在以上括号内填写课题或课题组负责人或实验室名称，未有此项声明内容的，可以不作特别声明）。

本人承诺辅修专业毕业论文（设计）（如有）的内容与主修专业不存在相同与相近情况。
}
\vspace{30mm}

\begin{flushright}
    {\xiaosi 学生声明（签名）：\hspace{6em}}\\[15mm]
    {\xiaosi 年\hspace{2em}月\hspace{2em}日\hspace{4em}}
\end{flushright}

\clearpage{\pagestyle{empty}\cleardoublepage}

% ==================== 致谢 ====================
\chapter*{致谢}
\thispagestyle{plain}
\addcontentsline{toc}{chapter}{致谢}

{\xiaosi\song
在此填写致谢内容。致谢语应以简短的文字对课题研究与论文撰写过程中曾直接给予帮助的人员（例如指导教师、答疑教师及其他人员）表示自己的谢意。
}

\vfill
\rightline{\kai 二〇二六年\quad 月于厦门大学}

\clearpage{\pagestyle{empty}\cleardoublepage}

% ==================== 中文摘要 ====================
\chapter*{摘要}
\thispagestyle{plain}
\addcontentsline{toc}{chapter}{摘要}

{\xiaosi\song
在此填写中文摘要内容。摘要应具有独立性和自含性，语言精炼、明确，高度概括论文内容，以400字左右为宜。
}
\newline

{\xiaosi\hei\bfseries 关键词}{\xiaosi\song ：关键词一；关键词二；关键词三；关键词四；关键词五}

\clearpage{\pagestyle{empty}\cleardoublepage}

% ==================== 英文摘要 ====================
\chapter*{Abstract}
\thispagestyle{plain}
\addcontentsline{toc}{chapter}{Abstract}

{\xiaosi
Fill in the English abstract here. The English abstract should be consistent with the Chinese abstract in content.
}
\newline

{\xiaosi\bfseries Keywords}{\xiaosi : Keyword One; Keyword Two; Keyword Three; Keyword Four; Keyword Five}

\clearpage{\pagestyle{empty}\cleardoublepage}

% ==================== 中文目录 ====================
\tableofcontents

\clearpage{\pagestyle{empty}\cleardoublepage}

% ==================== 英文目录 ====================
\tableofengcontents

\clearpage{\pagestyle{empty}\cleardoublepage}

% ============================================================
%                       正文部分
% ============================================================
\mainmatter

% ==================== 第一章 ====================
\chapter{绪论}
\echapter{Introduction}

在此填写绪论（引言）内容。引言应包括论文的研究目的、流程和方法等。研究领域的历史回顾、文献回溯、理论分析等内容应独立成章，用足够的文字叙述。

正文一般不少于6000字（不含图表、程序和计算数字）。

\section{研究背景}
\esection{Research Background}

在此填写研究背景内容。

\section{研究目的与意义}
\esection{Research Purpose and Significance}

在此填写研究目的与意义。

\section{论文结构安排}
\esection{Organization of the Thesis}

本文的结构安排如下：

第一章为绪论，介绍研究背景、研究目的与意义。

第二章为……

第三章为……

\clearpage{\pagestyle{empty}\cleardoublepage}

% ==================== 第二章 ====================
\chapter{相关理论基础}
\echapter{Theoretical Foundations}

在此填写相关理论基础内容。

\section{节标题示例}
\esection{Section Title Example}

这是正文内容示例，使用小四号宋体。以下展示各种排版元素的用法。

\subsection{小节标题示例}
\esubsection{Subsection Title Example}

这是小节内容示例。

% 公式示例
\subsubsection*{公式示例}

公式应另起一行居中书写，编号用英文圆括号括起，放在公式右边行末：
\begin{equation}
    E = mc^{2}
\end{equation}

多行公式示例：
\begin{align}
    f(x) &= a_{0} + a_{1}x + a_{2}x^{2} + \cdots \\
    g(x) &= b_{0} + b_{1}x + b_{2}x^{2} + \cdots
\end{align}

% 定理示例
\begin{thm}[示例定理]
    设 $X$ 和 $Y$ 是拓扑空间，$f: X \to Y$ 是连续映射。若 $X$ 是紧致的，则 $f(X)$ 也是紧致的。
\end{thm}

\begin{proof}
    在此填写证明过程。
\end{proof}

\begin{dfn}[示例定义]
    设 $V$ 是域 $F$ 上的向量空间，$T: V \to V$ 是线性变换……
\end{dfn}

% 图片示例
\subsubsection*{图片示例}

图号和图名要居图下方的正中（注意行首无缩进）。

\begin{figure}[htbp]
    \centering
    % \includegraphics[width=0.6\textwidth]{example-figure.png}
    \caption{示例图片标题}
    \label{fig:example}
\end{figure}

% 表格示例 - 三线表
\subsubsection*{表格示例}

表号和表名居表上方正中（注意行首无缩进）。表格应优先采用三线表。

\begin{table}[htbp]
    \centering
    \caption{数据集基本属性}
    \label{tab:example}
    \begin{tabular}{cccc}
        \toprule
        \textbf{数据集} & \textbf{规模} & \textbf{连续属性} & \textbf{离散属性} \\
        \midrule
        housing & 506 & 12 & 1 \\
        servo   & 167 & 0  & 4 \\
        \bottomrule
    \end{tabular}
\end{table}

\clearpage{\pagestyle{empty}\cleardoublepage}

% ==================== 第三章 ====================
\chapter{研究方法}
\echapter{Research Methods}

在此填写研究方法相关内容。

\clearpage{\pagestyle{empty}\cleardoublepage}

% ==================== 第四章 ====================
\chapter{实验结果与分析}
\echapter{Experimental Results and Analysis}

在此填写实验结果与分析内容。

\clearpage{\pagestyle{empty}\cleardoublepage}

% ==================== 第五章 ====================
\chapter{结论}
\echapter{Conclusion}

结论应包含论文的核心观点，阐述自己的创造性成果及其在本研究领域中的意义、作用，交代研究工作的局限，提出未来工作的意见和建议。

\clearpage{\pagestyle{empty}\cleardoublepage}

% ============================================================
%                     参考文献
% ============================================================
\backmatter

\chapter*{参考文献}
\addcontentsline{toc}{chapter}{参考文献}

% 参考文献格式应符合 GB/T 7714-2005
% 以下为手动编排示例，也可使用 BibTeX/BibLaTeX 管理

\begin{thebibliography}{99}

% 期刊文章 [J]
\bibitem{ref1}
作者1, 作者2, 作者3. 文章题目[J]. 期刊名, 出版年份, 卷号(期数): 起止页码.

% 专著 [M]
\bibitem{ref2}
作者. 书名(版次)[M]. 出版地: 出版单位, 出版年份: 起止页码.

% 学位论文 [D]
\bibitem{ref3}
作者. 论文题目[D]. 保存地: 保存单位, 年份.

% 英文期刊示例
\bibitem{ref4}
Author A, Author B. Title of the article [J]. Journal Name, 2024, 10(2): 100-120.

% 英文专著示例
\bibitem{ref5}
Author C. Book title [M]. City: Publisher, 2023.

\end{thebibliography}

\clearpage{\pagestyle{empty}\cleardoublepage}

% ============================================================
%                       附录
% ============================================================
\begin{appendices}
\renewcommand{\thechapter}{附录~\Alph{chapter}}

\chapter{附录标题示例}

在此填写附录内容。附录的篇幅不宜太多，一般不超过正文。

论文附录依次用大写字母"附录 A、附录 B、附录 C……"表示。

\end{appendices}

\end{document}
